%% ch01 --- Model to maszyna do przewidywania
%%
%% Napisany we wrześniu 2026 jako wzorzec dla wszystkich kolejnych rozdziałów.
%% Każda liczba, której czytelnik nie policzy w pamięci, pochodzi z
%% code/measure/ch01.py, który importuje TE SAME reguły, które drukują listingi
%% (code/ch01/), więc strona i listing nie mogą się rozejść.
%%
%% Bliźniak: chapters/en/ch01-models.tex. Podrozdział za podrozdziałem, makro
%% za makrem, liczba za liczbą; tools/parity.py je porównuje.

\chapter{Model to maszyna do przewidywania}
\label{ch:models}
\index{model}

\noindent
Korzystasz z modeli matematycznych codziennie, nawet jeśli tak ich nie
nazywasz. Aplikacja, która mówi, że pociąg przyjedzie za siedem minut, liczy
jakiś model. Tak samo minutnik w piekarniku, tablica pływów, dawkowanie na
ulotce leku i cichy głos, który podpowiada, że trzeba wyjść dziesięć po
ósmej, żeby zdążyć na dziewiątą. Każdy z nich bierze kilka liczb o tym, co
jest teraz, stosuje regułę i podaje ci liczbę o tym, co będzie później.

Ta książka opowiada o tym, co dzieje się z tą obietnicą. Pierwsze cztery
rozdziały pokazują, jak dobrze działa \dash{} jak linijka algebry potrafi
powiedzieć, gdzie będzie kamień, jak gorąca będzie twoja herbata i kiedy
wróci kometa. Reszta książki pokazuje, gdzie się łamie, dlaczego łamie się
nawet wtedy, gdy reguła jest doskonała, i co wciąż da się przewidzieć, gdy
to się stanie. To ostatnie jest powodem, by ją przeczytać: to pęknięcie nie
jest porażką matematyki, a znajomość jego kształtu zmienia to, w co
powinieneś wierzyć przy każdej prognozie, jaką ktoś ci poda.

\begin{outcomes}
\outcome{nazwać stan, regułę i przewidywanie każdego modelu, który ktoś ci
  pokaże}
\outcome{uruchomić model napisany w Pythonie i przeczytać to, co wypisuje}
\outcome{powiedzieć, co znaczy, że model jest błędny, i dlaczego błędny
  model wciąż może być użyteczny}
\outcome{powiedzieć, co dzieje się z małym błędem stanu początkowego w dwóch
  modelach z tego rozdziału}
\outcome{przytoczyć zegarmistrzowskie twierdzenie Laplace'a z 1814 roku i
  pytanie, które reszta książki mu zadaje}
\end{outcomes}

\section{Trzy części każdego modelu}
\label{sec:ch01-parts}
\index{model!stan}\index{model!reguła}\index{model!przewidywanie}

\noindent
Upuść kamień z mostu. Ile przeleci w ciągu trzech sekund?

Fizyk odpowie wzorem, który możesz pamiętać ze szkoły: kamień spadający
swobodnie przy powierzchni Ziemi pokonuje drogę
\[ d = \tfrac{1}{2}\, g\, t^{2} \]
w ciągu $t$ sekund, gdzie $g$ to mniej więcej \val{ch01.g} metra na sekundę
na sekundę. Ta jedna linijka to kompletny model i ma trzy części, które
znajdziesz w każdym modelu w tej książce.

\begin{itemize}
\item \textbf{Stan}: liczby, które mówią, jak jest teraz. Tutaj to jedna
  liczba \dash{} czas $t$ od chwili puszczenia kamienia.
\item \textbf{Reguła}: jak te liczby zamieniają się w to, co chcesz wiedzieć.
  Tutaj to wzór.
\item \textbf{Przewidywanie}: to, co mówi reguła. Tutaj \dash{} droga.
\end{itemize}

\begin{think}
Podstaw do wzoru $t = 3$. Mniej więcej ile przeleciał kamień: około $10$
metrów, około $30$, czy około $45$?
\end{think}
\reveal{Około \val{ch01.stone.three} metra: połowa z \val{ch01.g}, razy
$3 \times 3 = 9$.}
Jeśli odpowiedziałeś dziesięć albo trzydzieści, jesteś w dobrym
towarzystwie: większość ludzi zaniża, bo kamień nie spada ze stałą
prędkością. Ciągle przyspiesza, więc droga rośnie z \emph{kwadratem} czasu.
Trzy razy dłuższy lot to dziewięć razy dłuższa droga.

To pierwsza rzecz, którą daje ci model: łapie intuicję, która miała cię
zmylić. Nikt nie ma wyczucia dla $t^{2}$. Wzór go nie potrzebuje.

\begin{trapbox}
\textbf{\enquote{Model to kopia rzeczywistości.}} Nie jest, i nic w
powyższym wzorze nie przypomina kamienia. Model to maszyna do odpowiadania
na pytania: wkładasz liczby, dostajesz liczby, a jedyny uczciwy test polega
na tym, czy te liczby zgadzają się z tym, co potem zmierzysz. Mapa nie jest
terenem, a model nie jest nawet mapą \dash{} to mapa, która dodatkowo mówi
ci, o której dojedziesz.
\end{trapbox}

\section{Twoje laboratorium}
\label{sec:ch01-lab}
\index{Python}\index{laboratorium}

\noindent
Każdą liczbę w tej książce, której nie dało się policzyć w pamięci, policzył
mały program, a każdy z tych programów leży w repozytorium książki w
katalogu \code{code/}. Nie musisz być programistą, żeby je czytać. Są
krótkie, napisane tak, żeby dało się je czytać, i każdy jest objaśniony przy
pierwszym pojawieniu się. Oto kamień.

\pyregion{code/ch01/falling_stone.py}{model}{Spadający kamień jako model: jedna stała, jedna reguła}{lst:ch01-stone}

\noindent
Czytaj go tak, jak czytasz wzór. Wiersz zaczynający się od \code{G =}
ustala stałą. Wiersze zaczynające się od \code{def distance} definiują
regułę: podaj jej czas \code{t}, a zwróci \code{0.5 * G * t * t}, czyli wzór
z rozpisanymi mnożeniami. Tekst w potrójnych cudzysłowach to notatka dla
czytelnika, a wszystko po znaku \code{\#} to komentarz. Potem program
zadaje regule pytanie dla każdej pełnej sekundy od $0$ do $5$:

\pyregion{code/ch01/falling_stone.py}{table}{Sześć pytań zadanych modelowi}{lst:ch01-table}

\noindent
\code{for t in range(6)} znaczy \enquote{wykonaj następny wiersz z
\code{t} równym $0$, potem $1$ i tak dalej aż do $5$}. Uruchom go z katalogu
\code{code/} poleceniem \code{uv run python ch01/falling\_stone.py}, a
wypisze:

\transcript{ch01-falling-stone}

\noindent
Spójrz na prawą kolumnę. W pierwszej sekundzie kamień spada około pięciu
metrów, a w pięć sekund \val{ch01.stone.five} metra. Jeśli kiedykolwiek
upuściłeś coś z wysokości i liczyłeś, ta druga liczba wydaje się za duża,
a wcale taka nie jest.

\begin{note}
Materiał wstępny mówi, jak zainstalować laboratorium: jedno narzędzie,
\pkg{uv}, które instaluje dokładnie tę wersję Pythona, na której sprawdzono
książkę (\pypatch), i każdy pakiet w dokładnie tej wersji, którą podaje
strona redakcyjna. Potem \code{uv run python} z nazwą pliku uruchamia
dowolny listing z książki. Jeśli wolisz niczego nie uruchamiać, wynik
każdego listingu jest wydrukowany obok; uruchamianie ich to sposób, żeby
przekonać się, że książka nie prosi cię o wiarę na słowo.
\end{note}

\begin{lab}{l01_01_stone}{Kamień w obie strony}
Model odpowiada na pytania w więcej niż jednym kierunku. Podaj mu czas, a
poda drogę; podaj drogę, a poda czas; podaj zmierzony spadek, a powie, ile
musi wynosić grawitacja. Plik startowy zawiera trzy funkcje z opisami i bez
treści. Uzupełnij je, aż test przejdzie. Druga potrzebuje pierwiastka
kwadratowego, który w Pythonie nazywa się \code{math.sqrt}.
\end{lab}

\section{Model, który robi kroki}
\label{sec:ch01-steps}
\index{model!krokowy}\index{prawo stygnięcia Newtona}

\noindent
Wzór na kamień pozwala przeskoczyć od razu do dowolnej chwili. Większość
modeli na to nie pozwala. Nalej herbatę o temperaturze
\val{ch01.tea.start}\,°C w pokoju, w którym jest \val{ch01.tea.room}\,°C.
Isaac Newton opisał w 1701 roku, jak stygnie, a w postaci, której używa ta
książka, jego reguła mówi: \emph{co minutę herbata traci stały ułamek
różnicy między sobą a pokojem}. Niech ten ułamek wynosi \val{ch01.tea.k},
czyli jedną dziesiątą.

\begin{think}
Na początku różnica wynosi $90 - 20 = 70$ stopni. Jaka będzie temperatura
minutę później?
\end{think}
\reveal{\val{ch01.tea.one}\,°C. Jedna dziesiąta różnicy to $7$, a
$90 - 7 = 83$.}
Żeby znać temperaturę po dwóch minutach, robisz to samo jeszcze raz, z
$83$: różnica wynosi teraz $63$, jedna dziesiąta z niej to \num{6.3}, a
herbata ma \num{76.7}. Nie ma skrótu, który powinieneś znać. Reguła mówi, jak
przejść od jednej minuty do następnej, a żeby dojść gdziekolwiek indziej,
robisz krok po kroku.

\pyregion{code/ch01/cooling_cup.py}{rule}{Reguła następnej minuty}{lst:ch01-rule}

\noindent
Funkcja \code{next\_minute} to cały model. Dostaje temperaturę teraz i zwraca
temperaturę minutę później. Żeby zobaczyć przyszłość, wywołujesz ją raz za
razem, podając jej z powrotem każdą odpowiedź:

\pyregion{code/ch01/cooling_cup.py}{run}{Każda odpowiedź wraca na wejście, trzydzieści razy}{lst:ch01-run}

\transcript{ch01-cooling-cup}

\noindent
Wiersz \code{temp = next\_minute(temp)} to najważniejszy wiersz tej książki i
zobaczysz go w każdym rozdziale pod inną nazwą. Zastępuje stan następnym
stanem. Robienie tego raz za razem nazywa się \emph{iteracją}, a
rozdział~\ref{ch:iteration} jest wyłącznie o niej.

\begin{think}
Z wydruku wynika, że herbata ma więcej niż $60$\,°C w minucie $5$. Czy
ma mniej niż $60$ w minucie $6$? Co trzeba zrobić, żeby mieć pewność?
\end{think}
\reveal{Tak: pierwszy raz spada poniżej $60$ w minucie
\val{ch01.tea.drinkable}. Żeby mieć pewność, robisz regułą jeszcze jeden
krok albo prosisz program o każdą minutę zamiast co piątej.}
Tak zadaje się pytanie modelowi krokowemu: robisz kroki, aż pojawi się
odpowiedź. Dlatego też modele krokowe były rzadkie przed komputerami, a
teraz są wszędzie. Maszynie nie przeszkadza milion kroków.

\plotfig{ch01-two-models}{Dwa modele tego samego rodzaju pytania. Wzór na
kamień podaje dowolną chwilę od razu; reguła herbaty podaje następną minutę,
a krzywa to jej kroki połączone ze sobą.}{wzór i reguła krokowa}

\begin{lab}{l01_02_cooling}{Kiedy da się to wypić?}
Napisz \code{cool}, które stosuje regułę przez zadaną liczbę minut i
zapamiętuje każdą temperaturę, oraz \code{minutes\_until}, które odpowiada na
pytanie, jakie człowiek naprawdę zadaje. Twoja odpowiedź dla kubka o
$90$\,°C, pokoju o $20$ i ułamka \num{0.1} powinna się zgadzać z tą powyżej.
\end{lab}

\section{Błędny i użyteczny}
\label{sec:ch01-wrong}
\index{model!błąd}

\noindent
Żaden z tych modeli nie jest prawdziwy. Wzór na kamień pomija powietrze, a
piórko upuszczone obok kamienia pokazuje, jak bardzo może to mieć znaczenie.
Reguła herbaty pomija kubek, spodek, parę i twoją łyżeczkę. Zmierz dokładnie,
a przekonasz się, że oba są błędne.

\begin{think}
Skoro wzór na kamień pomija powietrze, a powietrze oczywiście ma znaczenie,
to czy ten wzór jest bezużyteczny?
\end{think}
\reveal{Nie. Dla kamienia spadającego kilka sekund jest dokładny co do
szerokości dłoni; dla piórka myli się bardzo. Model jest użyteczny do pytań,
na które odpowiada wystarczająco dobrze, a twoim zadaniem jest wiedzieć,
które to pytania.}

\begin{trapbox}
\textbf{\enquote{Model, który jest błędny, jest bezużyteczny.}} Każdy model
jest błędny, bo każdy coś pomija \dash{} właśnie dlatego jest dość mały, żeby
go używać. Często cytowane zdanie statystyka George'a Boxa mówi to samo
krócej: wszystkie modele są błędne, ale niektóre są użyteczne. Użyteczne
pytanie nigdy nie brzmi \enquote{czy jest prawdziwy?}, tylko \enquote{czy
jest wystarczająco dobry do tego pytania, w tym zakresie?}. Zapamiętaj to
pytanie. Część~II tej książki dotyczy zakresu, w którym model bez żadnego
błędu w regule przestaje być wystarczająco dobry, a powód nie jest tym,
czego byś się spodziewał.
\end{trapbox}

\noindent
Model można sprawdzić pomiarem, i to w obie strony. Upuść kamień z
dziesięciometrowego muru, a wzór powie, że wyląduje po
\val{ch01.stone.ten.time} sekundy. Zmierz prawdziwy spadek, dostań trochę
inną liczbę, a możesz odwrócić wzór i zapytać, ile musiałoby wynosić $g$. To
trzecia funkcja z pierwszego laboratorium powyżej i tak model zdobywa
zaufanie: formułując przewidywania na tyle ostre, że mogłyby okazać się
błędne \dash{} i nie okazując się.

\section{Co dzieje się z małym błędem?}
\label{sec:ch01-error}
\index{błąd!stanu początkowego}

\noindent
Oto pytanie, na którym opiera się cała książka, zadane najłagodniejszemu z
jej modeli. Nigdy nie znasz stanu początkowego dokładnie. Termometr pokazuje
$90$\,°C, a herbata ma naprawdę $91$. O ile pomylisz się pół godziny później?

\begin{think}
Dwa kubki zaczynają z różnicą jednego stopnia i stygną według tej samej
reguły. Po trzydziestu minutach różnica między nimi jest większa niż jeden
stopień, wciąż równa jednemu stopniowi, czy mniejsza?
\end{think}
\reveal{Mniejsza: około \val{ch01.gap.thirty}\,°C. Co minutę różnica jest
mnożona przez \val{ch01.gap.factor}, bo każdy kubek traci jedną dziesiątą
własnej odległości od temperatury pokoju.}

\pyregion{code/ch01/two_cups.py}{gap}{Dwa kubki różniące się o jeden stopień}{lst:ch01-gap}

\transcript{ch01-two-cups}

\noindent
Błąd nie tylko pozostaje mały; znika. Co minutę jest mnożony przez
\val{ch01.gap.factor}, więc po dziesięciu minutach wynosi mniej więcej
trzecią część stopnia, a po \val{ch01.gap.hundredth} minutach mniej niż
setną część. W skali logarytmicznej, w której każdy krok w dół to
dziesięciokrotność, różnica spada wzdłuż linii prostej:

\plotfig{ch01-gap-shrinks}{Różnica między dwoma kubkami, które zaczęły z
różnicą jednego stopnia. W skali logarytmicznej spada wzdłuż linii prostej:
co minutę znika ten sam ułamek.}{błąd, który zanika}

\noindent
Zapamiętaj tę prostą. Model, którego błędy maleją, wybacza ci, że nie znasz
stanu początkowego. Model kamienia jest prawie równie łaskawy: pomyl się o
jedną dziesiątą sekundy co do chwili puszczenia, a twoje przewidywanie
będzie przesunięte o stały czas na zawsze, nigdy o więcej. Większość
matematyki, którą poznałeś w szkole, opisuje takie układy i dlatego myśl,
że przyroda jest przewidywalna, wydaje się zdrowym rozsądkiem.

\begin{worldbox}
Dawkowanie na ulotce leku opiera się na modelu tego rodzaju. Większość leków
jest usuwana z krwi w mniej więcej stałym ułamku na godzinę, dokładnie tak,
jak herbata traci stały ułamek swojej różnicy (alkohol to słynny wyjątek),
więc dawka przyjęta trochę za późno albo trochę za duża zostawia błąd, który
maleje, a nie rośnie. Ta wyrozumiałość sprawia, że schemat dawkowania w ogóle
da się wydrukować na opakowaniu.
\end{worldbox}

\section{Mechanizm zegara}
\label{sec:ch01-clockwork}
\index{demon Laplace'a}\index{kometa Halleya}\index{wszechświat jak zegar}

\noindent
W XVIII wieku modele takie jak te dokonały czegoś zdumiewającego. Edmond
Halley, korzystając z Newtonowskich praw ruchu i grawitacji, dowodził, że
komety widziane w latach 1531, 1607 i 1682 to jeden obiekt na jednej
orbicie, i przewidział, że wróci około 1758 roku. Zmarł w 1742. Kometę
zobaczono ponownie w noc Bożego Narodzenia 1758 roku i od tamtej pory nosi
jego imię.

Reguła na papierze sięgnęła siedemdziesiąt sześć lat w przyszłość i miała
rację. Pierre-Simon Laplace wyciągnął z tego wniosek w 1814 roku, we
fragmencie, który w wolnym przekładzie brzmi:

\begin{quote}
Umysł, który w danej chwili znałby wszystkie siły ożywiające przyrodę i
wzajemne położenie wszystkich tworzących ją bytów \dash{} umysł dość
rozległy, by poddać te dane analizie \dash{} objąłby jednym wzorem ruchy
największych ciał wszechświata i najlżejszego atomu; nic nie byłoby dla
niego niepewne, a przyszłość, tak jak przeszłość, stałaby mu przed oczami.
\end{quote}

\noindent
Ten wyobrażony umysł nazywa się \emph{demonem Laplace'a}, a obraz, który za
nim stoi, \emph{wszechświatem jak mechanizm zegara}: znaj regułę i
teraźniejszość, a przyszłość wynika z nich tak pewnie jak ruch wskazówek
zegara. W 1960 roku fizyk Eugene Wigner pisał o \enquote{niepojętej
skuteczności matematyki w naukach przyrodniczych} i miał na myśli to samo
zdumienie: reguły spisane przez ludzi przy biurkach wciąż okazują się mieć
rację co do świata.

\mermaidfig{ch01-model-loop}{Każdy model w tej książce ma ten kształt. Stan
wchodzi, reguła zamienia go w przewidywanie, a pomiar decyduje, czy mu
ufać.}{stan, reguła, przewidywanie, sprawdzenie}

\begin{think}
Załóżmy, że reguła jest dokładna \dash{} nie w przybliżeniu poprawna, tylko
idealnie poprawna \dash{} i że ją znasz. Czy możesz wtedy przewidywać
przyszłość dowolnie daleko?
\end{think}
\reveal{Tylko jeśli dokładnie znasz też stan początkowy, a nigdy go nie
znasz. Dla herbaty i kamienia nie ma to prawie znaczenia, bo błąd stanu
początkowego maleje albo pozostaje taki sam. Czy zawsze zachowuje się tak
łaskawie \dash{} na to pytanie odpowiada reszta tej książki.}

\noindent
Większość ludzi odpowiada \enquote{tak}, i tak samo odpowiadała większość
naukowców przez dwieście lat po Laplace'u. Zapamiętaj swoją odpowiedź. W
rozdziale~\ref{ch:logistic} poznasz regułę długości jednej linijki, prostszą
niż reguła herbaty, której błędy wcale nie maleją; w
rozdziale~\ref{ch:butterfly} zmierzysz, jak szybko rosną; a pod koniec
książki będziesz umieć powiedzieć, liczbami, które sam policzysz, dlaczego
niektóre rzeczy da się przewidzieć na sto lat naprzód, a innych nawet na
dwa tygodnie. To jest teoria chaosu, i nie jest to teoria nieładu. To
matematyka dokładnie tego pytania.

\begin{summarybox}
\begin{itemize}
\item Każdy model ma trzy części: \textbf{stan} (liczby, które mówią, jak
  jest teraz), \textbf{regułę} (jak się zmieniają albo jak zamieniają się w
  odpowiedź) i \textbf{przewidywanie} (co mówi reguła).
\item Niektóre modele to wzory, które da się obliczyć dla dowolnej chwili.
  Większość to \textbf{reguły krokowe}: mówią, jak przejść od teraz do
  następnej chwili, a do przyszłości dochodzisz, stosując je raz za razem.
  To jest iteracja.
\item Każdy model jest błędny, bo każdy coś pomija. Użyteczne pytanie brzmi:
  czy jest wystarczająco dobry do tego pytania, w tym zakresie.
\item W przypadku stygnącej herbaty błąd stanu początkowego \textbf{maleje}
  co minutę o ten sam czynnik; w przypadku kamienia pozostaje taki sam.
  Takie modele wybaczają ci, że nie znasz stanu początkowego dokładnie.
\item Demon Laplace'a to zegarmistrzowskie twierdzenie: znaj dokładnie
  regułę i teraźniejszość, a przyszłość z nich wynika. Reszta książki je
  sprawdza.
\end{itemize}
\end{summarybox}

\canyou

\begin{checkpoint}
\item Nazwij stan, regułę i przewidywanie rozkładu jazdy autobusu.
  \answerto{Stanem jest aktualna godzina i przystanek, na którym jesteś;
  regułą jest rozkład (razem z założeniem, że autobus się go trzyma);
  przewidywaniem jest godzina przyjazdu następnego autobusu.}
\item Ile przeleci kamień w ciągu pierwszych dwóch sekund, bez oporu
  powietrza? A w ciągu czterech?
  \answerto{Około \num{19.6} metra w dwie sekundy i około \num{78.5} w
  cztery. Dwa razy dłuższy czas, cztery razy dłuższa droga.}
\item Herbata o $80$\,°C stygnie w pokoju o $20$\,°C, tracąc co minutę
  jedną dziesiątą różnicy. Jaką ma temperaturę po jednej minucie, a jaką po
  dwóch?
  \answerto{$74$\,°C, a potem \num{68.6}\,°C: różnica zmienia się z $60$ na
  $54$, a potem na \num{48.6}.}
\item Dwa kubki zaczynają z różnicą dwóch stopni w tym samym pokoju i stygną
  według tej reguły. O ile różnią się po jednej minucie?
  \answerto{O \num{1.8} stopnia: każdy kubek traci jedną dziesiątą własnej
  różnicy do temperatury pokoju, więc ich różnica też traci jedną
  dziesiątą.}
\item Dlaczego model, o którym wiadomo, że jest błędny, wciąż może być
  użyteczny?
  \answerto{Bo każdy model coś pomija, a użyteczność to kwestia tego, czy
  jego odpowiedzi są wystarczająco bliskie dla zadawanego pytania, w
  zakresie, w którym się go używa.}
\item Co musiałby znać demon Laplace'a, żeby przewidzieć przyszłość, i której
  z tych rzeczy według reszty książki nigdy nie da się znać dokładnie?
  \answerto{Regułę i dokładny stan obecny. Regułę można znać dokładnie; stanu
  obecnego nigdy, bo każdy pomiar ma błąd.}
\end{checkpoint}
